% !PRINTED 718
\weberpage{718}

Accordingly there results from \S.\,194, (12), having regard to the trigonometric formula
\[
\sin\frac{3\pi}{8} = \cos\frac{\pi}{8}
\]
\[
X_1 = -\frac{1}{\sqrt{2}} \log \operatorname{tg}\frac{\pi}{8}, \quad X_2 = \frac{-\pi}{4}, \quad X_3 = \frac{-\pi}{2\sqrt{2}}.
\]
Now
\[
\tau = \frac{r-1}{r^2(r+1)} = \operatorname{tg}\frac{\pi}{8} = \sqrt{2} - 1
\]
is a unit of the field $H_8$, which here is nothing other than the quadratic field arising from $\sqrt{2}$, and the units in it are therefore obtained from the solutions of the \gesperrt{Pell} equation
\[
u^2 - 2v^2 = \pm 1
\]
in the form $u + v\sqrt{2}$.

But the least positive solution of this \gesperrt{Pell} equation is $u = 1$, $v = 1$, and consequently one finds, by Bd.\,I, \S.\,128, all units in $H_8$ in the form $\pm(1+\sqrt{2})^a$, in which $a$ is a positive or negative integer exponent; $\tau$ results from this for $a = -1$, and hence all units are contained in the form
\[
\pm \tau^a
\]
as well. Thus $\tau$ is a \gesperrt{fundamental unit}, and the determinant $E$ occurring in the formulae of the preceding section [\S.\,195, (12)] here becomes (since $\tau$ is a proper fraction):
\[
E = -\log \tau.
\]
One thus obtains from \S.\,195, (15), (18):
\[
h_1 = 1, \quad k = 1, \quad h = 1.
\]
The field $\Omega_8$ is therefore a \gesperrt{one-class field}, i.e. one in which the decomposition of the integers into prime factors is possible according to the same laws as in the field of rational numbers.

At the same time it has emerged that all units in the field $H_8$ are representable in the form $\pm \tau^a$. But the conjugate values of the unit $\tau$ are
\[
\tau_1 = \operatorname{tg}\frac{\pi}{8}, \quad \tau_3 = -\operatorname{tg}\frac{3\pi}{8} = -\tau_1^{-1},
\]
and so have different signs. From this we conclude, for this case, the theorem:
%CONT>>
